\documentclass[10pt,a4paper]{article}
\usepackage[T1]{fontenc}
\usepackage[utf8]{inputenc}
\usepackage{lmodern}
\usepackage{amsmath,amssymb}
\usepackage[margin=21mm]{geometry}
\usepackage{microtype}
\usepackage[hidelinks]{hyperref}
\usepackage{enumitem}
\setlist[itemize]{leftmargin=*,nosep,topsep=3pt}
\setlength{\parindent}{0pt}
\setlength{\parskip}{5pt}
\setlength{\emergencystretch}{3em}
\newcommand{\question}[1]{\par\medskip{\normalsize\bfseries #1}\par\smallskip}
\pagestyle{plain}

\begin{document}
{\Large\bfseries Dopamine depletion leads to pathological synchronization of distinct basal ganglia loops in the beta band\par}

\textbf{Paper:} Ortone et al. (2023).\quad Spiking-network modelling study.\par
\href{https://doi.org/10.1371/journal.pcbi.1010645}{DOI: \nolinkurl{10.1371/journal.pcbi.1010645}}\par
{\footnotesize Page references below use PDF page numbers, including front matter.}

\question{To which general question does this paper relate?}
Possible network mechanisms of beta synchronization and the limitations of an isolated STN--GPe model.

\question{How does it answer the general questions?}
\textbf{Finding.} Interacting STN--GPe and striatal--GPe loops show beta-range resonances. In simulations, dopamine depletion strengthens the striatal loop and promotes synchronization between loops that are more independent at high dopamine levels (PDF pp.~1--4, Results).

\textbf{Usable insight.} The source of a shared spectral peak can depend on interactions between circuits and dopamine-dependent coupling. Network interventions should therefore be evaluated across mechanisms and operating conditions. The paper establishes a computational explanation of synchronization; it does not demonstrate that minimizing beta improves patients' motor function.

\textbf{Model relevance.} Six populations are included: D1/D2 medium spiny neurons, fast-spiking interneurons, two GPe populations and STN. Individual neurons have adaptation, conductance dynamics and reset events. These equations are not exactly the notebook's memoryless monotone rate model. Tables~2--3 give biophysical parameters and projection delays; they require a justified population reduction before becoming rate-model coefficients (PDF pp.~4--7).

\question{What papers does it cite that could provide more information?}
\begin{itemize}
\item De la Crompe et al. (2020), \emph{The globus pallidus orchestrates abnormal network dynamics in a model of Parkinsonism}: experimental constraint (ref.~49, PDF p.~28).
\item Corbit et al. (2016), \emph{Pallidostriatal projections promote beta oscillations in a dopamine-depleted biophysical network model}: striatal-loop comparison (ref.~57, PDF pp.~28--29).
\end{itemize}

\end{document}
